% ===================================================================
%  TÁBLA Uimh. 13 — An tÓstán. — An Bhialann. An Caife.
%  Traduction 2026 d'apres texto/io, controlee sur texto/fr, texto/en
%  et texto/es. Consigne : voir texto/ga/10-tabla-01.tex.
%
%  DEUX NOTES D'AUTEUR, l'une et l'autre en gras-italique dans le
%  fac-simile, comme dans l'ido : la premiere renvoie au tableau 6
%  pour le detail de la chambre, la seconde au vocabulaire pour la
%  carte des plats.
%
%  DEUX INVERSIONS, ET TOUTES DEUX VIENNENT ENCORE DU VERBE — la
%  huitieme et la neuvieme de cette colonne, toutes du meme cote.
%
%  LA SECONDE EST CELLE DU NOM VERBAL, pour la troisieme fois de la
%  colonne apres les tableaux 7 et 8. « Donacir al garsono (57) bel
%  gratifiketo (59) » : l'ido pose d'abord le destinataire, l'objet
%  ensuite. La tournure irlandaise « X a thabhairt do Y » place son
%  objet AVANT le nom verbal, donc avant le destinataire, et (59)
%  passait devant (57). Le remede est de quitter le nom verbal pour un
%  verbe conjugue, qui rend l'ordre verbe-sujet-destinataire-objet :
%  « thug mé don fhreastalaí (57) síneadh láimhe breá (59) ». La
%  regle est constante : le nom verbal irlandais antepose son objet,
%  le verbe conjugue le postpose.
%
%  LA PREMIERE EST CELLE DU COMPLEMENT DE MANIERE. L'ido ecrit
%  « mashinisto (52) reparis, segun l'indiki dil ciklisto (53),
%  petrol-triciklo (51) » : le complement de maniere s'intercale
%  entre le verbe et son objet. L'irlandais ne l'intercale pas — il
%  range verbe, sujet, objet, et rejette le complement a la fin, ce
%  qui faisait passer (51) devant (53). Le remede est celui de
%  toujours : refaire la phrase. On donne au mecanicien une relative
%  qui porte le complement — « meicneoir (52), a bhí ag obair de réir
%  threoracha an rothaí (53), i mbun deisiúcháin ar thríchruthach
%  peitril (51) » — et l'ordre de l'ido revient.
%
%  LE VOCABULAIRE DE L'HOTEL EST CELUI DE LA TABLE, et il se partage
%  comme au tableau 6 : ce qui se mange est irlandais — « im »,
%  « iasc », « cáis », « torthaí », « uachtar » — et ce qui se sert
%  arrive avec le service : « biachlár », « bialann », « meanú »,
%  « ardaitheoir », « teileagram », « ceamara », « gluaisteán ».
% ===================================================================

%%K t13-noto-1 noto
\VUnotes{143.2mm}{%
\textit{\VUgras{(*) I dtaobh mionsonraí an tseomra, féach an tábla Uimh. 6.}}%
}

%%K t13-apar-1 apar
\VUsaut{3.0mm}
\VUtitre{15.0pt}{60}{TÁBLA Uimh. 13}
\VUsaut{1.6mm}
\VUpk{10.2pt}{40}{An tÓstán. --- An Bhialann.}
\VUsaut{2.0mm}
\VUpk{10.2pt}{40}{An Caife.}
\VUsaut{3.4mm}

%%K t13-01-1 p
1. --- Ó tháinig mé go Royan tráthnóna inné, chuir mé fúm in « \textit{Óstán an
Aigéin} », áit a mhol cara dom.

%%K t13-01-2 p
Tugadh \VUgras{seomra} breá dom \textsuperscript{(2)} ar an dara \VUgras{hurlár}
\textsuperscript{(1)}, agus chodail mé go han-mhaith ann (*).

%%K t13-02-1 p
2. --- Ar maidin inniu, agus mé ag fágáil mo sheomra, áit ar thug an \VUgras{cailín
freastail} \textsuperscript{(3)} an bricfeasta chugam, chuaigh mé isteach sa
\VUgras{seomra tobac} \textsuperscript{(5)} a úsáidtear mar \VUgras{sheomra
léitheoireachta} \textsuperscript{(4)} chomh maith, chun na \VUgras{nuachtáin}
\textsuperscript{(6)} a léamh agus \VUgras{toitín} \textsuperscript{(8)} a chaitheamh.
Tar éis na léitheoireachta chuaigh mé síos san \VUgras{ardaitheoir}
\textsuperscript{(9)} agus amach ag siúl tríd an mbaile.

%%K t13-03-1 p
3. --- Ó rinne mé dearmad a fhiafraí den \VUgras{chléireach} \textsuperscript{(12)}
san óstán cén t-am a chuirtear na béilí ar an \VUgras{mbord coiteann}
\textsuperscript{(11)}, d'fhill mé go luath agus bhain mé leas as sin chun dul á
fholcadh féin i \VUgras{seomra folctha} \textsuperscript{(10)} an óstáin. Ansin
chuaigh mé arís chuig \VUgras{oifig an airgid} \textsuperscript{(15)} chun glaoch a
chur ar dhuine de mo chairde. Ach bhí an \VUgras{teileafón} \textsuperscript{(13)}
gafa; nuair a chonaic \VUgras{bean an airgid} \textsuperscript{(14)} mo chás, í ag
breacadh \VUgras{bille} \textsuperscript{(17)} ar \VUgras{chlár na dtaistealaithe}
\textsuperscript{(16)}, d'iarr sí ar an \VUgras{doirseoir} \textsuperscript{(18)} an
\VUgras{gasúr} \textsuperscript{(19)} a chur ag déanamh mo theachtaireachta \VUgras{ar
rothar} \textsuperscript{(20)}.

%%K t13-04-1 p
4. --- Ghabh mé buíochas léi agus chuaigh mé amach chun síneadh láimhe a thabhairt
don \VUgras{iompróir} \textsuperscript{(24)} (fear na hoibre trom) a thug ón stáisiún
ar a \VUgras{thralaí láimhe} \textsuperscript{(25)} mo \VUgras{bhosca hata}
\textsuperscript{(26)}, \VUgras{mo chás taistil} \textsuperscript{(27)} agus
\VUgras{mo mhála taistil} \textsuperscript{(28)}. Bhí \VUgras{fear an phoist}
\textsuperscript{(21)} tar éis teacht, agus thug sé \VUgras{litir}
\textsuperscript{(22)} dom.

%%K t13-05-1 p
5. --- D'fhan mé tamall eile ar leac an dorais, in aice leis an \VUgras{mbosca
litreach} \textsuperscript{(23)}, ag féachaint ar an siúl ar an tsráid. Bhí
\VUgras{tiománaí} \textsuperscript{(30)} an \VUgras{omnibus} \textsuperscript{(29)} ag
luchtú \VUgras{chófra} \textsuperscript{(31)} \VUgras{taistealaí}
\textsuperscript{(32)} ar a charr, taistealaí a raibh \VUgras{fear an óstáin}
\textsuperscript{(33)} á thionlacan. Thug \VUgras{teileagrafaí} óg
\textsuperscript{(35)} \VUgras{teileagram} \textsuperscript{(34)} isteach gan deifir
do chliant den óstán; bhí \VUgras{turasóir} \textsuperscript{(36)} agus a
\VUgras{cheamara} \textsuperscript{(38)} ar iompar ar a ghualainn aige, ag dul i
gcomhairle lena \VUgras{threoirleabhar} \textsuperscript{(37)}.

%%K t13-tit-1 sub
\VUsaut{4.0mm}
\VUpk{10.2pt}{40}{Am an Lóin.}
\VUsaut{3.0mm}

%%K t13-06-1 p
6. --- Os comhair na ráillí bhí \VUgras{teachtaire} \textsuperscript{(39)}
neamhchúiseach ina thámhnéal idir a \VUgras{chrúca iompair} \textsuperscript{(40)}
agus a \VUgras{bhosca snasáin} \textsuperscript{(41)}, fad is a bhí \VUgras{bean
níocháin} óg \textsuperscript{(42)} a bhí ag iompar \VUgras{ciseáin}
\textsuperscript{(43)} éadaigh ag gáire faoi ghnúis shollúnta \VUgras{mhaor an
bhoird} \textsuperscript{(44)} a bhí ina sheasamh in aice leis an doras.

%%K t13-07-1 p
7. --- Nuair a chonaic mé an \VUgras{cócaire} \textsuperscript{(45)}, a tháinig
amach as a \VUgras{chistin} \textsuperscript{(47)} atá san \VUgras{íoslach}
\textsuperscript{(48)} chun \VUgras{buachaill cistine} \textsuperscript{(46)} a chur
ar theachtaireacht, chuimhnigh mé go raibh uair an lóin ag druidim liom, agus chuaigh
mé isteach sa bhialann, ag trasnú na cúirte inar chuir \VUgras{tiománaí
gluaisteáin} \textsuperscript{(50)} a \VUgras{ghluaisteán} \textsuperscript{(49)} ar
leataobh, agus ina raibh \VUgras{meicneoir} \textsuperscript{(52)}, a bhí ag obair de
réir threoracha \VUgras{an rothaí} \textsuperscript{(53)}, i mbun deisiúcháin ar
\VUgras{thríchruthach peitril} \textsuperscript{(51)} a raibh timpiste bainte dó.

%%K t13-08-1 p
8. --- D'fhiafraigh mé de \VUgras{ghasúr} óg \textsuperscript{(54)} a bhí ag iompar
\VUgras{gloiní beorach} \textsuperscript{(55)}, an raibh an bord coiteann faoin
vearanda, agus ó dúirt sé go raibh, shuigh mé ag ceann de na boird agus thug mé orm
béile den scoth a fháil.

%%K t13-noto-2 noto
\VUnotes{143.2mm}{%
\textit{\VUgras{(*) Le cabhair na liosta miasa atá curtha sa fhoclóir,
féadfaidh an múinteoir an meanú thíos a athrú go héasca.}}%
}

%%K t13-tit-2 sub
\VUsaut{4.0mm}
\VUpk{10.2pt}{40}{Lón den Scoth.}
\VUsaut{3.0mm}

%%K t13-09-1 p
9. --- Thug an freastalaí \VUgras{biachlár} an lóin (*) chugam agus liosta na
bhfíonta. Roghnaigh mé agus d'ordaigh mé mar \VUgras{chéadchúrsa ribí róibéis} le
\VUgras{him,} agus, mar bhéile tosaigh, \VUgras{putóga ar nós Caen.} Níor ith mé
\VUgras{iasc,} ach d'iarr mé cuid de \VUgras{cheathrú} caoireola, agus, mar
\VUgras{mhias ghlasraí, spionáiste} le huachtar. Ansin, ó nár thaitin ceann ar bith
de na \VUgras{hidirchúrsaí} liom, chuir mé fios ar mhilseoga éagsúla,
\VUgras{cáis,}
\VUgras{torthaí} agus \VUgras{brioscaí.} Chuir mé leis sin \VUgras{fíon} den scoth ó
Saint-Émilion, agus, agus mé an-sásta le mo bhéile, thug mé don
\VUgras{fhreastalaí} \textsuperscript{(57)} \VUgras{síneadh láimhe} breá
\textsuperscript{(59)} agus d'fhág mé an bord.

%%K t13-10-1 p
10. --- Nuair a chonaic \VUgras{an bainisteoir} \textsuperscript{(60)} go raibh mé
réidh le himeacht, mhol sé dom dul chun an bhalcóin chun radharc iontach
\VUgras{an Aigéin} \textsuperscript{(62)} a fheiceáil. I gcéin thángthas ar
\VUgras{bháidíní} iascaigh \textsuperscript{(63)}, ar \VUgras{longa seoil}
\textsuperscript{(64)} agus ar \VUgras{phaicéad} ollmhór \textsuperscript{(65)} a
raibh a scáil dhubh ag seasamh amach ar \VUgras{scamaill} bhána
\textsuperscript{(67)} \VUgras{bhun na spéire} \textsuperscript{(66)}. Chuir leoithne
éadrom \VUgras{an bhratach} \textsuperscript{(70)} a bhí sáite os cionn
\VUgras{chomhartha} \textsuperscript{(69)} an \VUgras{halla} \textsuperscript{(68)} ag
luascadh.

%%K t13-tit-3 sub
\VUsaut{3.0mm}
\VUpk{10.2pt}{40}{Caitheamh Aimsire.}
\VUsaut{3.0mm}

%%K t13-11-1 p
11. --- Bhí an radharc chomh suimiúil sin gur shocraigh mé dul suas amárach ar
theirís an chaife agus \VUgras{déghloine} \textsuperscript{(71)} nó \VUgras{gloine
fhada} \textsuperscript{(72)} a bhreith liom.

%%K t13-12-1 p
12. --- Ag fágáil an bhalcóin dom chuaigh mé chuig caife an óstáin chun
\VUgras{deoch} \textsuperscript{(73)} a ól agus \VUgras{cluiche billéardaí}
\textsuperscript{(75)} a imirt. Ghabh mé gan stad tríd an seomra caife inar imir a
lán custaiméirí \VUgras{ficheall} \textsuperscript{(78)}, \VUgras{táiplis}
\textsuperscript{(79)}, \VUgras{dúradáin} \textsuperscript{(80)}, \VUgras{táiplis
mhór} \textsuperscript{(81)} nó \VUgras{cártaí} \textsuperscript{(82)}, agus chuaigh
mé díreach suas chun an \VUgras{tseomra bhilléardaí} \textsuperscript{(74)}.

%%K t13-13-1 p
13. --- Nuair a bhí an cluiche thart, chuaigh mé síos go hurlár na talún agus d'iarr
mé ar an bhfreastalaí \VUgras{clúdach} \textsuperscript{(84)}, \VUgras{páipéar}
\textsuperscript{(83)}, \VUgras{stampa poist} \textsuperscript{(85)}, \VUgras{peann}
\textsuperscript{(87)} agus \VUgras{dúch} \textsuperscript{(86)} chun litir a
scríobh.

%%K t13-13-2 p
Ansin ghlaoigh mé ar \VUgras{chóisteoir} \textsuperscript{(92)} a raibh a
\VUgras{charráiste} \textsuperscript{(88)} stoptha os comhair an chaife. D'fhiafraigh
mé de faoin bpraghas de réir na n-uaireanta nó de réir an turais, agus, nuair a
thángamar ar aon fhocal faoin bpraghas, d'oscail sé \VUgras{doras an charráiste}
\textsuperscript{(89)} dom agus chuir sé isteach mé. Chuaigh sé suas arís ar an
\VUgras{suíochán} \textsuperscript{(91)}, chuir sé na capaill chun siúil go mear leis
an \VUgras{bhfuip} \textsuperscript{(93)}, agus d'imíomar ar shiúlóid álainn.
